\begin{table}[ht!]
\centering
\footnotesize
\caption{The as-of backtest: what the register's lag alone produces in years where the true age gap is observable.}
\label{tab:iv_backtest}
\begin{tabular}{lccc}
\toprule
Register truncated at & True codes & As-of codes & Artefact \\
\midrule
2021 & $+0.0193$ (0.0129) & $-0.2875$ (0.0171) & $\mathbf{-0.3068}$ \\
2022 & $+0.0176$ (0.0111) & $-0.1452$ (0.0122) & $\mathbf{-0.1627}$ \\
\bottomrule
\end{tabular}
\begin{minipage}{0.90\textwidth}\footnotesize\vspace{4pt}The coefficient is the submitted design's $\hat\gamma_2$ on 2019--2023, years in which every worker's true contemporaneous occupation is observed. The as-of column re-runs it having truncated the register at the stated year and rebuilt the coding cascade as it would have been, which imposes on those years exactly the staleness that 2024--25 inherit. \textbf{The artefact is the gap between the two columns, not the as-of coefficient.} At the 2021 truncation the lag moves an estimate of $+0.019$, indistinguishable from zero, to $-0.288$; the artefact is $-0.307$. A design of that kind returns about $-0.17$ on these data, so the lag alone can produce more than the whole of it. The design the paper reports admits no occupation code recorded after 2019, so this test does not apply to it. Its own 2019 codes are not all fresh: restricting the score to those observed in 2019, rather than carried back from earlier files, moves its adoption step by $+0.005$, and 97.8 per cent of employers keep their quartile. Source: scripts 45, 68 and 82.\end{minipage}
\end{table}
